\documentclass[11pt,a4paper]{article}

\usepackage[margin=1in]{geometry}
\usepackage[T1]{fontenc}
\usepackage[utf8]{inputenc}
\usepackage{lmodern}
\usepackage{setspace}
\usepackage{graphicx}
\usepackage{booktabs}
\usepackage{hyperref}
\usepackage{enumitem}
\usepackage{float}

\onehalfspacing
\title{MindGuard: Emotion-Aware and Crisis-Sensitive NLP Prototype}
\author{Team Name \\ Module: Advance AI}
\date{\today}

\begin{document}
\maketitle

\begin{abstract}
This report presents MindGuard, an academic prototype that combines multi-label emotion classification, crisis-risk detection, and empathetic response generation. We evaluate model quality, safety behavior, and operational reliability under realistic constraints.
\end{abstract}

\section{Introduction}
\subsection{Problem Statement}
\subsection{Objectives and Scope}
\subsection{Contributions}

\section{Related Work}
Summarize relevant literature and compare methods used in this project.

\section{Data and Preprocessing}
\subsection{Datasets}
\subsection{Cleaning and Normalization}
\subsection{Splits and Versioning}

\section{Methodology}
\subsection{System Architecture}
\subsection{Emotion Classification with BERT}
\subsection{Crisis Detection with VAE}
\subsection{Response Generation Strategy}

\section{Implementation}
\subsection{Tech Stack}
\subsection{API and UI Design}
\subsection{MLOps and Experiment Tracking}

\section{Evaluation and Results}
\subsection{Emotion Metrics}
\subsection{Crisis Metrics and Baseline Comparison}
\subsection{Response Quality Assessment}

\section{Responsible AI, Ethics, and Limitations}
\subsection{Safety Policy}
\subsection{Bias and Fairness Risks}
\subsection{Non-Clinical Usage Disclaimer}

\section{Conclusion and Future Work}

\bibliographystyle{plain}
\bibliography{references}

\appendix
\section{Reproducibility Checklist}
Include environment, commands, model artifacts, and run IDs.

\end{document}
